\documentclass[11pt]{article}
\usepackage[a4paper,margin=25mm]{geometry}
\usepackage{amsmath,amssymb,amsthm,booktabs}
\newtheorem{proposition}{Proposition}
\title{The two-colour cubic example as an affine Coxeter geometry}
\author{}\date{}
\begin{document}\maketitle
Example 2.5 of Fernandes--Leemans--Weiss uses red cubes, yellow cubes,
edges and vertices of the ordinary cubic tessellation. We identify its
actual incidence structure with the universal coloured-vertex model
underlying the acquired halving construction. This argument uses the
infinite tessellation directly; it does not set a finite period parameter
equal to zero outside the hypotheses of a toroid theorem.

\paragraph{A coloured-vertex geometry.}
In the unit cubic tessellation of $\mathbb R^3$, take objects of types
0 and 1 to be the even and odd vertices of $\mathbb Z^3$, respectively.
Types 2 and 3 are its squares and cubes. Opposite-colour vertices are
incident when they differ by a signed coordinate vector. Incidence
with a face, or between faces, is containment. Put
\[
D=\{v\in\mathbb Z^3:v_1+v_2+v_3\equiv0\pmod2\},
\qquad K=D\rtimes B_3,
\]
where $B_3$ is the full signed permutation group.

\begin{proposition}
This incidence structure is a regular hypertope. Its full
type-preserving group is $K$. With the base chamber
\[
 0,\quad e_1,\quad [0,1]^2\times\{0\},\quad [0,1]^3,
\]
its distinguished involutions are
\[
 x=\tau_{e_1+e_2}z_1z_2s_1,\qquad s_1,\qquad s_2,\qquad z_3.
\]
Here $s_i$ exchanges coordinates $i,i+1$, $z_i$ reverses coordinate
$i$, and composition acts from the rightmost factor first. The diagram
has bonds $02$ and $12$ labelled 3, and bond $23$ labelled 4.
\end{proposition}
\begin{proof}
A chamber is an edge, with its two differently coloured endpoints,
contained in a square contained in a cube. Every flag extends to such
a chain. Indeed its existing squares and cubes are nested; any specified
vertices lie in its smallest face. If both vertices occur they form
an edge of that face, since they are unit-step neighbours. Extend that
edge, or a specified vertex, through the missing face dimensions.
With no prescribed face one may choose an incident unit square and cube.

There are exactly two completions of any corank-one flag. For a missing
vertex these are the two neighbours of the other vertex in its square.
For a missing square they are the two faces of the given cube containing
the specified edge. For a missing cube they are the two cubes containing
the given square. Thus the geometry is thin.

Its two-type residue incidence graphs are cycles with the following
lengths:
\[
\begin{array}{c|rrrrrr}
\text{remaining types}&01&02&12&03&13&23\\\hline
\text{cycle length}&4&6&6&4&4&8.
\end{array}
\]
For $02$ and $12$, the residue records the three edges and three
faces at a corner of a fixed cube. For $23$, it records the four
squares and four cubes around a fixed edge. The other cases are the
four corners of a square, or the two relevant corners together with
the two cubes on that square.

Every three-type residue is connected as well. At a cube the ordinary
cube graph connects all its vertex objects, and every square meets it.
At a square the four corners and its two containing cubes form a
connected graph. At a vertex, its six neighbours are connected through
the squares: every two neighbours on different coordinate axes belong
to a square through that vertex. Each incident cube meets this connected
part. The whole structure is connected because the unit grid connects
all vertex objects and every face contains a vertex. This proves
residual connectivity, and the two-type calculation gives the diagram.

Translations in $D$ move any even vertex to the origin. Signed coordinate
permutations then act simply transitively on the 48 chains from that
vertex: six choices of edge, four of containing square and two of cube.
Hence $K$ acts regularly on all chambers. The four displayed maps fix
exactly the other three objects of the base chamber, and move their own
object to its other completion. The acquired universal affine calculation
proves that these maps generate $D\rtimes B_3$ with precisely the affine
Coxeter presentation stated above. In particular the chamber graph is
connected. A type-preserving automorphism fixing a chamber fixes every
adjacent chamber by thinness, hence every chamber. Thus there are no
additional type-preserving automorphisms.
\end{proof}

\newpage
\paragraph{The literal identification with the source.}
Colour a source cube $v+[0,1]^3$ red when $\sum v_i$ is even, and yellow
otherwise. Use the dual cubic grid with vertices
$\frac12\mathbf1+\mathbb Z^3$. A source cube maps to its centre.
A source vertex $v$ maps to the dual cube
$v+[-\frac12,\frac12]^3$. A source edge
$v+[0,1]e_i$ maps to the dual square whose $i$th coordinate is
$v_i+\frac12$, and whose other coordinates range from
$v_j-\frac12$ to $v_j+\frac12$.

These maps are bijections onto the respective dual cells. A source
vertex is in an edge or cube precisely when the corresponding dual
square or vertex is contained in its dual cube. Likewise an edge lies
in a source cube precisely when its dual square contains that cube's
centre. Finally two source cubes of different colours share a square
precisely when their centres are unit-step neighbours. This checks
every pair of distinct types, not just chambers.

Translate the dual grid by $-\frac12\mathbf1$. Source types
$(0,2,1,3)$ then become types $(0,1,2,3)$ of the proposition:
red cubes become even vertices, yellow cubes odd vertices, edges
squares and vertices cubes. Thus the source example is exactly this
regular hypertope, with the indicated permutation of types. In the
source numbering the bonds $01$ and $21$ have label 3 and bond $13$
has label 4, agreeing with its native diagram.

\paragraph{Acquired theory, checks and attribution.}
The equality and Coxeter presentation
$\langle x,s_1,s_2,z_3\rangle=D\rtimes B_3$ are the universal affine
lemma in the accepted generalized-halving calculation. They precede
all finite period quotients. The coloured-vertex model is also the
universal model underlying that construction. The present addition is
the explicit incidence bijection and the proof of the infinite example
at the exact interface required by this source.

Small integer-affine checks verify the four displayed transformations,
their exact pair orders, their three fixed base objects and their
distinct fourth object. The 48 signed coordinate maps give 48 distinct
base-vertex chambers. These checks accompany the direct proof; finite
sampling is not used to prove residual connectivity of an infinite grid.
The example and its red/yellow interpretation belong to
Fernandes--Leemans--Weiss. The earlier halving construction and affine
calculation retain their stated sources and authorship. This note proves
one source example and asserts no new classification, conjecture result,
whole-manuscript disposition, formal verification or compression rate.
\end{document}
